
\documentclass[10pt,a4]{article}
\usepackage[english]{babel}
\usepackage[utf8x]{inputenc}
\usepackage{amsmath}
\usepackage{graphicx}
\usepackage[colorinlistoftodos]{todonotes}
\usepackage{hyperref}
\usepackage{geometry}
\geometry{top=3cm,left=2cm,right=2cm,bottom=3cm}
\usepackage[scaled]{helvet}
\usepackage[T1]{fontenc}
\renewcommand\familydefault{\sfdefault}


\author{Samuele Lo Piccolo}
\date{\today}
\title{Driver for Adafruit Neo Trinkey - SAMD21 USB Key}



\begin{document}
\maketitle
\tableofcontents

\section{Project data}

\begin{itemize}
\item 
  Project supervisor: Vittorio Zaccaria

\item 
The group that is delivering this project:

\begin{center}
\begin{tabular}{lll}
Last and first name & Person code & Email address\\
\hline
  Lo Piccolo Samuele & 10796320 & samuele.lopiccolo@mail.polimi.it \\               
\end{tabular}
\end{center}

\item The Github repository can be found \href{https://github.com/Samuele-LP/AOS}{\textbf{here}}

\end{itemize}


\section{Project description}



This projects allows to communicate with a serial interface exposed as a USB-CDC endpoint on the Adafruit Neo Trinkey.\\
The aim of the project is to be able to send textual commands to the Circuit Python REPL and control the device's behaviour, mainly to change the LEDs' colours.\\\\
I did not implement the handling of multiple devices since I only have access to one Neo Trinkey and could not \\reliably validate the correctness of the solution


\subsection{Design and implementation}

My driver registers itslef with the USB subsytem as a USB driver, providing only the required fields of \textbf{struct usb\textunderscore driver}:
\begin{itemize}
	\item The driver name
	\item The USB device table to be able to call the probe function
	\item The probe function that readies the driver and the device for communication
	\item The disconnect function that handles the disconnection of the device
\end{itemize}
The physical device is exposed as a character device in the /dev/ directory, the only file operations implemented are read and write.
\subsubsection{Init}
The Init function registers my driver to USBcore so that when a match for the \textit{Module Device Table} is found \hyperref[sec:probe]{\textbf{Probe}} will be called.\\
Init then goes on to: \begin{itemize}
	\item Request a character device region 
	\item Initializing and registering with the kernel the character device associated with the module
	\item Creating a class of devices "aos\textunderscore class"
\end{itemize}
\subsubsection{Exit}
The Exit function starts by unregistering from USBcore so that the disconnect function is called if a device is currently plugged in, then it deletes and unregisters the character device and finally it destroys the created class.
\subsubsection{Probe}
\label{sec:probe}
The Probe function is used to set the necessary CDC parameters for communication, it does so by sending two \textit{USB Control Messages} to the device.\\The first contains the \textbf{Line Coding} parameters which determine how the device and the host communicate.\\The second control message instead sets up the CDC endpoint to accept bulk transfers, in particular I found that setting both DTR and RTS to 1 worked best.\\
If the device setup is successful I create a device file \textbf{/dev/AOSdev}
\subsubsection{Disconnect}
When Disconnect is called I destroy the device file \textbf{/dev/AOSdev} and empty a local buffer of characters read from the device, but not yet sent to userspace.
\subsubsection{Write}
In the Write function I make use of USB bulk transfers to send data to the device, since \textbf{usb\textunderscore bulk\textunderscore msg} only accepts kernel memory I must first create a local buffer and copy the userspace buffer to it.
The device accepts up to 64 bytes at a time, so any write bigger than 64 bytes is truncated and the number of actually written bytes is returned to the caller.\\\\
The device has a local buffer of, I think, 64 bytes that contains the textual outputs of the REPL interaction, if this buffer ever fills up the device stops accepting bulk transfers; in the Write function I make sure that the device buffer is always empty after a write by performing successive USB bulk transfer reads until no more bytes are read, the output of these reads is then stored on a local buffer called \textit{read\textunderscore buffer} that is used when the \hyperref[sec:read]{\textbf{Read}} function is called.
\subsubsection{Read}
\label{sec:read}
The Read function does not communicate directly with the physical device, it just copies to userspace the contents of \textit{read\textunderscore buffer}\\
My buffer is of fixed size and gets flushed when it overflows or after a full read is completed \textbf{i.e.} when the number of read bytes is greater or equal to the number of available bytes. 

\section{Project outcomes}

\subsection{Concrete outcomes}
The produced artifacts are available in \href{https://github.com/Samuele-LP/AOS}{\textbf{Github}} and are:
\begin{itemize}
	\item This report
	\item The project presentation
	\item The device driver
\end{itemize}

\subsection{Learning outcomes}
\begin{itemize}
\item This project helped me to understand how to work more effectively in a Linux environment
\item I learned how USB endpoint enumeration works and the basics of how the Linux kernel implements communication with USB devices
\end{itemize}

\subsection{Existing knowledge}
I think most of the relevant knowledge was given by AOS, the rest of my career did not cover these topics with this level of depth if at all.\\
A course that helped me better understand C programming in general was Computer Security.

\subsection{Problems encountered}
The main issue of the project was understanding how my device worked, I could not find any documentation so I had to figure out why it behaved like that by trial and error. This is my first project using a physical device, so this is my first time tackling this type of problems.\\

\section{Honor Pledge}
(\textbf{This part cannot be modified and it is mandatory to sign it})

I/We pledge that this work was fully and wholly completed within the criteria
established for academic integrity by Politecnico di Milano (Code of Ethics and
Conduct) and represents my/our original production, unless otherwise cited.

I/We also understand that this project, if successfully graded,  will fulfill part B requirement of the
Advanced Operating System course and that it will be considered valid up until
the AOS exam of Sept. 2026. 

\begin{flushright}
Samuele Lo Piccolo
\end{flushright}


\end{document}
